%%%%%%%%%%%%%%%%%%%%%%%%%%%%%%%%%%%%%%%%%%%%%%%%%%%%%%%%%%%
\begin{frame}[fragile]\frametitle{}
\begin{center}
{\Large AI for Project and Program Managers}\\[0.5em]
{\large Unpack Curiosity Beyond the Job Title}
\end{center}
\end{frame}

%%%%%%%%%%%%%%%%%%%%%%%%%%%%%%%%%%%%%%%%%%%%%%%%%%%%%%%%%%%
\begin{frame}[fragile]\frametitle{}
\begin{center}
{\Large Why This Matters Now}

{\small AI and the PM Profession}
\end{center}
\end{frame}

%%%%%%%%%%%%%%%%%%%%%%%%%%%%%%%%%%%%%%%%%%%%%%%%%%%%%%%%%%%
\begin{frame}[fragile]\frametitle{Why PM and Program Manager Roles Are Facing Disruption}

\begin{itemize}
  \item A large share of a PM's week is spent on tasks AI is now good at:
    \begin{itemize}
      \item Drafting status reports, schedules, and meeting minutes
      \item Breaking a scope statement into tasks (Work Breakdown Structure)
      \item Scanning updates across a dozen threads to spot a risk before it becomes an issue
    \end{itemize}
  \item These are exactly the tasks Generative and Agentic AI are strongest at: reading, structuring, summarizing, and pattern-spotting across large amounts of text.
  \item \textbf{The uncomfortable question:} if AI can draft the plan, spot the risk, and summarize the meeting, what is left for the PM to do?
\end{itemize}

\end{frame}

%%%%%%%%%%%%%%%%%%%%%%%%%%%%%%%%%%%%%%%%%%%%%%%%%%%%%%%%%%%
\begin{frame}[fragile]\frametitle{Myths vs.\ Reality, What Every PM Must Know}

\begin{itemize}
  \item \textbf{Myth:} AI will replace the Project Manager.\\
    \textbf{Reality:} AI replaces the \textit{mechanical} parts of the job, drafting, formatting, first-pass analysis, not the judgment, negotiation, and accountability parts.

  \item \textbf{Myth:} AI plans are ready to execute as-is.\\
    \textbf{Reality:} AI drafts a plausible-looking plan from limited context. It does not know your vendor's real reliability, your sponsor's politics, or your team's actual velocity. Every AI output needs a PM's judgment applied before it becomes a real plan.

  \item \textbf{Myth:} Using AI is a shortcut that signals you're not doing the job properly.\\
    \textbf{Reality:} The PMs who use AI well are not doing \textit{less} work, they're doing \textit{higher-value} work: more time on stakeholders, negotiation, and unblocking, less time re-typing the same status update five ways.

  \item \textbf{Myth:} You need to learn to code to use any of this.\\
    \textbf{Reality:} Everything in this session runs on plain-English prompts, no coding required.
\end{itemize}

\end{frame}

%%%%%%%%%%%%%%%%%%%%%%%%%%%%%%%%%%%%%%%%%%%%%%%%%%%%%%%%%%%
\begin{frame}[fragile]\frametitle{The Hinton Lesson, A Role Is Not One Task}

\begin{itemize}
  \item In 2016, Geoffrey Hinton predicted radiologists would be obsolete within five years.
    A decade on, radiologists are still here, and in demand.
  \item \textbf{Why?} Radiology is not one task, it's roughly 24--25 distinct tasks. AI did
    automate one of them, reading an image for a pattern, better than most humans. It did not
    automate the rest: weighing an unusual finding against a patient's history, spotting an
    unrelated comorbidity, deciding what to tell an anxious patient.
  \item \textbf{The PM parallel:} your role is also a bundle of tasks, not one task. Some
    (drafting, summarizing, first-pass scheduling) will be automated. Others (judgment,
    negotiation, accountability) won't.
  \item \textbf{The real skill going forward:} learning to tell which of your tasks is which.
\end{itemize}

\end{frame}

%%%%%%%%%%%%%%%%%%%%%%%%%%%%%%%%%%%%%%%%%%%%%%%%%%%%%%%%%%%
\begin{frame}[fragile]\frametitle{From Task Execution to Strategic Orchestration}

\textbf{The PM's new value proposition:}
\begin{itemize}
  \item \textbf{Old value:} I can build the schedule. I can track the risks. I can write the status report.
  \item \textbf{New value:} I can direct AI to produce a first-draft schedule, risk register, and status report in minutes, then apply the judgment, context, and stakeholder relationships that make it \textit{actually usable}.
  \item \textbf{Analogy:} A PM using AI well is less like a clerk filling in templates and more like an editor-in-chief: setting direction, reviewing drafts critically, and deciding what goes out the door.
\end{itemize}

\vspace{0.3em}
\textbf{Plain English:} AI is the world's fastest, tirelessly patient junior project coordinator. It still needs a PM to tell it what the real project is trying to achieve, and to catch what it gets wrong.

\end{frame}

%%%%%%%%%%%%%%%%%%%%%%%%%%%%%%%%%%%%%%%%%%%%%%%%%%%%%%%%%%%
\begin{frame}[fragile]\frametitle{}
\begin{center}
{\Large AI-Powered Project Planning}
\end{center}
\end{frame}

%%%%%%%%%%%%%%%%%%%%%%%%%%%%%%%%%%%%%%%%%%%%%%%%%%%%%%%%%%%
\begin{frame}[fragile]\frametitle{What Is AI-Powered Project Planning?}

\begin{itemize}
  \item \textbf{Traditional approach:} PM manually decomposes scope into a Work Breakdown Structure (WBS), estimates each task, sequences them, and builds a schedule, often starting from a blank template or last project's plan.
  \item \textbf{AI-assisted approach:} PM describes the project's goal, scope, and constraints in plain English; AI proposes a first-draft WBS, task list, and rough sequencing in seconds.
  \item \textbf{What doesn't change:} the PM still validates every task against real team capacity, real dependencies, and real organizational constraints AI cannot see.
\end{itemize}

\end{frame}

%%%%%%%%%%%%%%%%%%%%%%%%%%%%%%%%%%%%%%%%%%%%%%%%%%%%%%%%%%%
\begin{frame}[fragile]\frametitle{From Manual WBS to AI-Assisted WBS}

\begin{center}
\begin{tabular}{|p{2.1cm}|p{2.4cm}|p{2.4cm}|}
\hline
\textbf{Step} & \textbf{Manual (today)} & \textbf{AI-Assisted (now available)} \\
\hline
First draft & 1--2 hours from a blank template & Under 5 minutes from a project brief \\
Task granularity & Depends on PM's memory of past projects & Consistent, pulls patterns from thousands of similar plans \\
Estimation & PM's gut feel or historical spreadsheet & AI proposes a starting estimate; PM adjusts using real team data \\
Missed tasks & Common, especially in unfamiliar domains & AI flags standard tasks a PM may not have worked with before \\
Ownership of the plan & PM & Still PM: AI drafts, PM decides \\
\hline
\end{tabular}
\end{center}

\end{frame}

%%%%%%%%%%%%%%%%%%%%%%%%%%%%%%%%%%%%%%%%%%%%%%%%%%%%%%%%%%%
\begin{frame}[fragile]\frametitle{Hands-On Exercise: Generating a Work Breakdown Structure}

\begin{itemize}
  \item \textbf{Activity:} Turn a one-line project idea into a structured WBS.
  \item \textbf{Try this prompt} (with any general-purpose AI assistant):\\
    \textit{``You are an experienced Project Manager. I am launching [your project, one sentence]. Break this into a 3-level Work Breakdown Structure: Phases, Deliverables, and Tasks. For each task, suggest a rough duration and the skill needed.''}
  \item \textbf{Discuss:} What did it get right? What did it miss that only you would know?
  \item \textbf{Key insight:} AI gives you a strong first draft in minutes. Your job shifts from \textit{creating} the WBS to \textit{critiquing and correcting} it.
\end{itemize}

\end{frame}

%%%%%%%%%%%%%%%%%%%%%%%%%%%%%%%%%%%%%%%%%%%%%%%%%%%%%%%%%%%
\begin{frame}[fragile]\frametitle{Prompting Patterns for Project Planning}

\begin{itemize}
  \item \textbf{Role Prompting:} ``You are a senior Program Manager in [your industry]. Review this schedule for unrealistic durations.''\\
    \textit{Why it works: Activates domain-appropriate scrutiny, not generic advice.}

  \item \textbf{Context + Constraints:} Always give AI: the project goal, team size, hard deadlines, and any constraints (budget, vendor dependency, regulatory approval).\\
    \textit{``Given a 6-person team and a fixed launch date of [date], which of these 20 tasks are at risk of slipping?''}

  \item \textbf{Structured Output:} Ask for a table, a numbered list, or a specific format (Gantt-ready CSV, RAID log columns), so the output plugs directly into your existing tools.

  \item \textbf{Iterate:} Treat the first AI output as a draft. ``Now re-sequence assuming Task 4 depends on an external vendor.''
\end{itemize}

\end{frame}

%%%%%%%%%%%%%%%%%%%%%%%%%%%%%%%%%%%%%%%%%%%%%%%%%%%%%%%%%%%
\begin{frame}[fragile]\frametitle{}
\begin{center}
{\Large AI Decision Support}
\end{center}
\end{frame}

%%%%%%%%%%%%%%%%%%%%%%%%%%%%%%%%%%%%%%%%%%%%%%%%%%%%%%%%%%%
\begin{frame}[fragile]\frametitle{What Is AI Decision Support?}

\begin{itemize}
  \item Not AI \textit{making} the decision, AI \textit{laying out the trade-offs} clearly and quickly, so the PM decides faster and with more complete information.
  \item Common PM decisions AI can help structure:
    \begin{itemize}
      \item Slip the deadline, add resources, or cut scope, which trade-off best serves the goal?
      \item Which of three vendors best fits this project's risk profile?
      \item Given current velocity, is the committed date still realistic?
    \end{itemize}
  \item \textbf{The PM's role doesn't shrink here, it sharpens:} deciding \textit{which} trade-off matters most to this sponsor, this team, this moment, is still entirely a human judgment call.
\end{itemize}

\end{frame}

%%%%%%%%%%%%%%%%%%%%%%%%%%%%%%%%%%%%%%%%%%%%%%%%%%%%%%%%%%%
\begin{frame}[fragile]\frametitle{Decision Support Use Cases}

\begin{center}
\begin{tabular}{|p{2.5cm}|p{3.8cm}|}
\hline
\textbf{Decision} & \textbf{How AI Supports It} \\
\hline
Schedule vs.\ scope trade-off & Lists what each option costs/saves, in time, budget, and risk \\
Resource reallocation & Simulates ``what if we move 2 people from Task A to Task B'' \\
Vendor/tool selection & Summarizes trade-offs across 3--4 options against your stated criteria \\
Go/no-go on a milestone & Synthesizes status data into a single readiness score with reasoning shown \\
\hline
\end{tabular}
\end{center}

\vspace{0.3em}
\textbf{Critical rule:} Always ask AI to show \textit{why}, not just \textit{what}. A recommendation without visible reasoning is not decision support, it's a guess dressed up confidently.

\end{frame}

%%%%%%%%%%%%%%%%%%%%%%%%%%%%%%%%%%%%%%%%%%%%%%%%%%%%%%%%%%%
\begin{frame}[fragile]\frametitle{Hands-On Exercise: AI Decision-Support Prompt}

\begin{itemize}
  \item \textbf{Activity:} Use AI to lay out a real trade-off you're facing (or a sample one).
  \item \textbf{Try this prompt:}\\
    \textit{``I am 3 weeks behind on a 12-week project. I can: (a) add 2 contractors for the last 4 weeks, (b) cut a non-critical feature, or (c) push the launch date by 2 weeks. For each option, list the cost, the risk, and who I'd need to convince.''}
  \item \textbf{Discuss:} Did AI surface an option or risk you hadn't considered? Where would you push back on its reasoning?
  \item \textbf{Key insight:} This is not about AI choosing for you, it's about walking into the sponsor conversation with all three options already priced out.
\end{itemize}

\end{frame}

%%%%%%%%%%%%%%%%%%%%%%%%%%%%%%%%%%%%%%%%%%%%%%%%%%%%%%%%%%%
\begin{frame}[fragile]\frametitle{}
\begin{center}
{\Large Live Demo}

{\small From Project Brief to Plan}
\end{center}
\end{frame}

%%%%%%%%%%%%%%%%%%%%%%%%%%%%%%%%%%%%%%%%%%%%%%%%%%%%%%%%%%%
\begin{frame}[fragile]\frametitle{Live Demo Setup}

\textbf{Sample project brief} (used for the rest of this demo):

\vspace{0.3em}
\textit{``Launch a mobile app for internal expense reporting at a 200-person company. Timeline: 10 weeks. Team: 1 PM, 2 developers, 1 designer, 1 QA. Must integrate with the existing HR system. Depends on an external vendor for the payment-gateway module.''}

\vspace{0.5em}
\begin{itemize}
  \item Follow along on your own device with any general-purpose AI assistant.
  \item We'll build, live, in front of the room: a WBS, a RAID log, and one real decision-support answer, from this single brief.
\end{itemize}

\end{frame}

%%%%%%%%%%%%%%%%%%%%%%%%%%%%%%%%%%%%%%%%%%%%%%%%%%%%%%%%%%%
\begin{frame}[fragile]\frametitle{Step 1: AI Generates a Work Breakdown Structure}

\textbf{Prompt used:}\\
\textit{``Using the project brief above, generate a 3-level WBS: Phases, Deliverables, Tasks. Flag any task that depends on the external payment-gateway vendor.''}

\vspace{0.3em}
\textbf{What to watch for in the live output:}
\begin{itemize}
  \item Does it correctly isolate the vendor-dependent tasks?
  \item Are the phase names sensible for a 10-week timeline, or overly generic?
  \item What would \textit{you} add that AI, lacking context on this specific company's HR system, couldn't know?
\end{itemize}

\end{frame}

%%%%%%%%%%%%%%%%%%%%%%%%%%%%%%%%%%%%%%%%%%%%%%%%%%%%%%%%%%%
\begin{frame}[fragile]\frametitle{Step 2: AI Generates a Risk \& Dependency Register}

\textbf{Prompt used:}\\
\textit{``From the same brief, generate a RAID log. Pay special attention to the HR system integration and the external payment-gateway vendor as dependency sources.''}

\vspace{0.3em}
\textbf{What to watch for in the live output:}
\begin{itemize}
  \item Did it flag the vendor dependency as the highest-impact risk (it should)?
  \item Did it surface any risk category you hadn't listed yourself (e.g. HR-system data-access approvals)?
  \item Notice: the log is a strong \textit{starting point} for the room's discussion, not a finished risk assessment.
\end{itemize}

\end{frame}

%%%%%%%%%%%%%%%%%%%%%%%%%%%%%%%%%%%%%%%%%%%%%%%%%%%%%%%%%%%
\begin{frame}[fragile]\frametitle{Step 3: AI Decision Support, ``Should We Compress the Schedule?''}

\textbf{Prompt used:}\\
\textit{``The vendor says the payment-gateway module will be 2 weeks late. I can: (a) launch without payment integration and add it later, (b) delay the whole launch by 2 weeks, or (c) run QA in parallel with the vendor's remaining work. Lay out the trade-offs for each.''}

\vspace{0.3em}
\textbf{What to watch for in the live output:}
\begin{itemize}
  \item Are the trade-offs specific to \textit{this} project, or generic boilerplate?
  \item Which option would you actually recommend to the sponsor, and why might your answer differ from AI's framing?
\end{itemize}

\end{frame}

%%%%%%%%%%%%%%%%%%%%%%%%%%%%%%%%%%%%%%%%%%%%%%%%%%%%%%%%%%%
\begin{frame}[fragile]\frametitle{Debrief: What This Demo Shows, and Its Limits}

\begin{itemize}
  \item \textbf{What it shows:} in under 10 minutes, AI produced a first-draft WBS, a RAID log, and a structured decision analysis, work that would traditionally take a PM an hour or more to draft from scratch.
  \item \textbf{Its limits:} every output here was generic to the \textit{stated} brief. It knows nothing about your actual vendor's history, your team's real morale, or your sponsor's actual risk appetite.
  \item \textbf{The real workflow:} AI drafts in minutes, you refine with the context only you have, in minutes more. The total cycle is still far faster than starting from a blank page.
\end{itemize}

\end{frame}

%%%%%%%%%%%%%%%%%%%%%%%%%%%%%%%%%%%%%%%%%%%%%%%%%%%%%%%%%%%
\begin{frame}[fragile]\frametitle{}
\begin{center}
{\Large Practical Tools \& Workflows}

{\small You Can Adopt Starting Today}
\end{center}
\end{frame}

%%%%%%%%%%%%%%%%%%%%%%%%%%%%%%%%%%%%%%%%%%%%%%%%%%%%%%%%%%%
\begin{frame}[fragile]\frametitle{AI Tool Landscape for Project/Program Managers}

\begin{center}
\begin{tabular}{|p{2.5cm}|p{4.4cm}|}
\hline
\textbf{Tool Category} & \textbf{Examples \& What They're For} \\
\hline
General-purpose assistants & ChatGPT, Claude, Gemini, for WBS drafting, RAID logs, decision framing, status-report writing \\
Office-embedded AI & Microsoft Copilot (Teams meeting summaries, Excel trend analysis, Outlook drafting) \\
PM-tool-native AI & Jira/Azure DevOps AI features (auto-summarizing tickets, sprint-risk flags), Asana/Monday AI (status rollups) \\
AI agents \& automation & Tools that watch a project's data (tickets, calendars, docs) and proactively flag risks, still an emerging category, evaluate carefully \\
\hline
\end{tabular}
\end{center}

\vspace{0.3em}
\textbf{Key principle:} Start with the general-purpose assistant you already have. You don't need a specialized PM-AI product to get 80\% of today's value.

\end{frame}

%%%%%%%%%%%%%%%%%%%%%%%%%%%%%%%%%%%%%%%%%%%%%%%%%%%%%%%%%%%
\begin{frame}[fragile]\frametitle{Beyond One Prompt, Agentic Workflows for PMs}

\begin{itemize}
  \item A single prompt gives you one draft from one perspective. An \textbf{agentic workflow} chains
    several AI agents together, each with a narrow job, so the output is checked before it reaches you.
  \item \textbf{No-code/low-code route:} platforms like n8n, Zapier, or Make let you wire this up
    by connecting boxes, no programming required.
  \item \textbf{If you're a bit tech-savvy:} the same pattern can be built with agents, skills, or
    commands in a coding assistant like Claude Code.
  \item \textbf{Worked example, from one project brief:}
    \begin{itemize}
      \item Agent 1 drafts the WBS
      \item Agent 2 drafts the RAID log
      \item Agent 3 critiques both drafts against the brief
      \item Agent 4 synthesizes a decision-ready summary for you to review
    \end{itemize}
  \item \textbf{Still true here:} you review and own the final output, agents remove drafting effort, not accountability.
\end{itemize}

\end{frame}

%%%%%%%%%%%%%%%%%%%%%%%%%%%%%%%%%%%%%%%%%%%%%%%%%%%%%%%%%%%
\begin{frame}[fragile]\frametitle{}
\begin{center}
{\Large Staying Relevant}

{\small What To Learn Next}
\end{center}
\end{frame}

%%%%%%%%%%%%%%%%%%%%%%%%%%%%%%%%%%%%%%%%%%%%%%%%%%%%%%%%%%%
\begin{frame}[fragile]\frametitle{What To Learn Next}

\begin{itemize}
  \item \textbf{Prompting itself:} the single highest-leverage skill from today, sharper prompts mean better first drafts everywhere below.
  \item \textbf{Microsoft Copilot:} if your org already runs Microsoft 365, learn its Teams meeting summaries, Excel analysis, and Outlook drafting, no new procurement needed.
  \item \textbf{One no-code/low-code agent platform:} pick one of n8n, Zapier, or Make; build one small chained automation (WBS draft $\rightarrow$ RAID draft $\rightarrow$ summary) to feel how agentic workflows work.
  \item \textbf{If tech-savvy:} learn agents, skills, and commands in a coding assistant like Claude Code, same chaining idea, with more control.
\end{itemize}

\end{frame}

%%%%%%%%%%%%%%%%%%%%%%%%%%%%%%%%%%%%%%%%%%%%%%%%%%%%%%%%%%%
\begin{frame}[fragile]\frametitle{References}

\begin{itemize}
  \item Lets Learn Prompt Engineering, Yogesh Haribhau Kulkarni (companion cheat sheet covering the prompting patterns used throughout this session)
  \item Geoffrey Hinton, 2016 remarks predicting the end of radiologist training within five years
\end{itemize}

\end{frame}

%%%%%%%%%%%%%%%%%%%%%%%%%%%%%%%%%%%%%%%%%%%%%%%%%%%%%%%%%%%
\begin{frame}[fragile]\frametitle{Thank You}

\begin{itemize}
  \item Thank you for your curiosity and hands-on energy today!
  \item One question to carry back: ``What is the \textit{one} AI-assisted task I will try on my current project tomorrow?''
\end{itemize}

\vspace{0.5em}
\begin{center}
\textit{``The Project Managers who thrive in the AI era won't be the ones who type the fastest, they'll be the ones who ask AI the sharpest questions, then apply the judgment it doesn't have.''}
\end{center}

\end{frame}
